\Conclusion

В результате выполнения выпускной квалификационной работы достигнута поставленная цель: разработан проект клиент‑серверной системы корпоративного обучения с поддержкой адаптивного тестирования для персонализации учебного процесса сотрудников и студентов.

Выполнены сформулированные задачи: проанализированы предметная область и требования к системе, подходы к адаптивному тестированию и мультиарендной архитектуре; проведён анализ аналогов и выбор технологий; спроектированы функциональность, роли пользователей, архитектура и модель данных; описаны сценарии использования, алгоритм адаптивного подбора вопросов и изоляции данных по тенантам; представлены прототипы интерфейсов мобильного приложения и веб‑панели администратора. Подготовлены рекомендации по внедрению, тестированию и вопросам информационной безопасности.

Практическая значимость работы заключается в возможности использования полученных проектных решений и прототипов как основы для развёртывания корпоративной или вузовской LMS с адаптивным тестированием и мультиарендной поддержкой. В перспективе планируется реализация серверной части и клиентских приложений в соответствии с разработанной спецификацией требований, проведение испытаний и при необходимости апробация в учебном процессе или представление результатов на конференции.
